\chapter{Écrire des règles en trois fonctions}

Un moteur de règles complet tient en trois méthodes.

\begin{center}
\begin{tabular}{@{}p{3.6cm}p{9.8cm}@{}}
\toprule
\textbf{Méthode} & \textbf{La question à laquelle elle répond} \\
\midrule
\texttt{Construire} & à quoi la partie ressemble au départ \\
\texttt{CoupsPossibles} & ce qu'on a le droit de faire \\
\texttt{Appliquer} & ce qui se passe quand on le fait \\
\bottomrule
\end{tabular}
\end{center}

Tout le reste --- créer et détruire une instance, cloner un état, le sérialiser,
le hacher, dire si la partie est finie, dire qui a gagné, vérifier qu'un coup est
légal --- est \textbf{le même pour tout le monde}, et il est déjà écrit.

Le fichier de ce chapitre existe, il compile, et l'atelier le charge :
\nkfile{Applications/ConquerorLab/exemples/rules/RegleFacile.cpp}.

\section{Le compte, mesuré}

\texttt{NkcRulesVTable} compte \textbf{24 entrées}. Voici où elles vont.

\begin{nkcode}{ou vont les 24 entrees de la vtable}
24  entrees de la vtable
 3  contiennent VOTRE jeu    Construire, CoupsPossibles, Appliquer
18  sont ecrites une fois pour toutes dans ConquerorRegleFacile.h
 3  sont optionnelles (geometrie) et restent nulles
\end{nkcode}

Les recopier à chaque module, ce serait dix-huit occasions de se tromper sur du
code qui n'a aucun rapport avec le jeu qu'on essaie de concevoir.

\section{Le module entier}

Voici un moteur qui joue vraiment. Rien n'a été coupé.

\begin{nkcode}{exemples/rules/RegleFacile.cpp --- le module complet}
#include "Conqueror/ConquerorRegleFacile.h"

using namespace nkentseu;
using namespace nkentseu::conqueror;
using namespace nkentseu::conqueror::facile;

struct MesRegles {

    static constexpr int32 kCote = 5;

    // 1. A quoi la partie ressemble au depart.
    void Construire(Grille &g) {
        g.topologie = NkcTopology::Square4;
        g.nbJoueurs = 2;
        g.AjouterRectangle(kCote, kCote);
        g.PoserAuxCoins(2);
    }

    // 2. Ce qu'on a le droit de faire : dupliquer vers une case voisine VIDE.
    void CoupsPossibles(const Partie &p, ListeCoups &out) {
        NkcCoord voisins[8];
        for (int32 i = 0; i < p.nbCases; ++i) {
            if (p.cases[i].owner != static_cast<int8>(p.joueur)) continue;
            const NkcCoord de = p.ou[i];
            const int32    n  = p.Voisins(de, voisins, 8);
            for (int32 k = 0; k < n; ++k)
                if (p.Vide(voisins[k]))
                    out.Dupliquer(p.joueur, de, voisins[k]);
        }
        // PASSER est ajoute tout seul si cette liste reste vide.
    }

    // 3. Ce qui se passe quand on le fait.
    void Appliquer(Partie &p, const NkcMove &m, Evenements &ev) {
        p.Poser(m.to, m.player);
        ev.Duplique(m.player, m.from, m.to);
        p.conquete[m.player] += 10;   // en DIXIEMES : 10 == 1,0 point

        NkcCoord    voisins[8];
        const int32 n = p.Voisins(m.to, voisins, 8);
        for (int32 k = 0; k < n; ++k) {
            if (!p.Ennemie(voisins[k], m.player)) continue;
            const int8 ancien = p.Proprietaire(voisins[k]);
            p.Poser(voisins[k], m.player);
            ev.Retourne(m.player, voisins[k], ancien);
        }
        p.PasserLaMain();
    }
};

NKC_REGLES(MesRegles, "MesRegles", "1.0.0", "Moi")
\end{nkcode}

La dernière ligne fabrique les vingt et une entrées de la vtable et les deux
symboles exportés. \textbf{C'est tout le module.}

\section{L'essayer}

\begin{enumerate}
  \item copiez \texttt{exemples/rules/RegleFacile.cpp} vers
        \texttt{travail/rules/mes\_regles.cpp} ;
  \item changez le nom dans \texttt{NKC\_REGLES} --- sinon deux entrées
        identiques au menu ;
  \item sauvegardez, attendez une seconde ;
  \item panneau \textbf{Modules} $\rightarrow$ votre module apparaît
        $\rightarrow$ \textbf{Nouvelle partie}.
\end{enumerate}

Si la compilation échoue, la sortie complète du compilateur s'affiche dans le
panneau \textbf{Modules}. C'est votre seul retour : lisez-la.

\section{Ce que l'échafaudage vous coûte}

Il n'y a pas de magie, et la contrepartie se dit en une phrase : \textbf{\texttt{Partie}
est une structure fixe, plate, sans pointeur.}

C'est ce qui permet aux dix-huit fonctions d'être justes \textbf{par
construction} et non par relecture.

\begin{center}
\begin{tabular}{@{}p{5.2cm}p{8.2cm}@{}}
\toprule
\textbf{ce que le cadre doit faire} & \textbf{comment, parce que \texttt{Partie} est plate} \\
\midrule
cloner un état & une affectation \\
sérialiser & un \texttt{memcpy} \\
hacher & \textsc{fnv-1a} sur les octets \\
dire si un coup est légal & énumérer les coups possibles et comparer \\
\bottomrule
\end{tabular}
\end{center}

Le prix : vous ne pouvez pas ranger dans \texttt{Partie} un état de taille libre
--- une liste qui grandit, un cache, un arbre. Tant que votre jeu tient dans les
cases, les joueurs, l'énergie et la conquête, vous n'y penserez jamais.

\section{Quand partir --- et ce n'est pas un échec}

\begin{quote}
\textbf{C'est un échafaudage, pas une cage.}
\end{quote}

Le jour où votre état ne rentre plus dans \texttt{Partie}, vous écrivez le
contrat nu. Les vingt-quatre entrées sont documentées au chapitre
\textbf{«~Quand l'échafaudage ne suffit plus~»}, à la fin de ce cours --- et il
se lira beaucoup mieux à ce moment-là, parce que vous aurez déjà vu les dix-huit
fonctions à l'œuvre.

Vous n'avez rien à défaire pour passer de l'un à l'autre : les deux produisent le
même module, avec les mêmes deux symboles. \nkfile{RegleFacile.cpp} et
\nkfile{RegleContratNu.cpp} jouent \textbf{exactement le même jeu} ---
comparez-les, c'est l'exercice le plus instructif du cours.

\section*{Exercices}

\begin{enumerate}
  \item \textbf{Portée 2.} Autorisez la duplication vers une case à deux pas.
        Une seule ligne change dans \texttt{CoupsPossibles} --- laquelle ?
  \item \textbf{Le déplacement.} Ajoutez l'action \textsc{deplacer} (la case de
        départ se vide). Que devient le décompte des totems ?
  \item \textbf{Hexagone.} Passez \texttt{g.topologie} à
        \texttt{NkcTopology::HexPointy}. Combien de voisins avez-vous
        maintenant, et quelle ligne de votre code l'a supposé ?
  \item \textbf{La cascade.} Émettez \texttt{ev.Cascade(...)} quand un coup
        retourne au moins deux ennemis, et regardez ce que l'atelier en fait à
        l'écran.
\end{enumerate}
